% supp_checks.tex -- Appendix S4: the audits of the main verifier and the cross-checks in detail; known minor
% issues of the tools (Sections 7.5, 7.6 and the end of Section 7.7 of the first revision)
\section{The main verifier, audits and cross-checks: details}\label{sec:supp-checks}

This appendix states the leaf test of the main verifier step by step, with its statistics
(\cref{ssec:algleaf}); it gives the details of the audits of the main verifier and of the cross-checks,
which \cref{ssec:audits,ssec:cross} summarize; and it lists the known minor issues of the tools mentioned
in \cref{ssec:trust}. The description of the leaf test documents the program that is named first in
\cref{fact:H}; the variant of \cref{alg:leaf} that this program uses is justified in the main part
(\cref{lem:deleted}). None of the audits and cross-checks reported here is part of the proof of
\cref{thm:main}.

\subsection{The leaf test of the main verifier: steps and cost}\label{ssec:supp-mainleaf}

For the record, and because \cref{ssec:trust} asks the reader to trust an implementation, we state what
\fn{verify/treecheck.cpp} does after it has checked (L1) and computed a pairing, when $\theta\ge1$. All
steps are justified by \cref{cor:pruning}, \cref{rem:true-cert} and \cref{lem:deleted}.
\begin{enumerate}[label=(\arabic*),leftmargin=2.2em]
\item \emph{Quick rejection.} If some $F$-connected component $S$ of the live lines has at least three
  lines and $\sur_F(S)<\theta$, the verdict is false; this is a violation of (L3) by definition.
\item \emph{Single lines.} One backward search from the unassigned cells finds the lines from which an
  augmenting path exists with respect to the pairing; the other lines have gain $0$ and are weak. A line
  $l$ with $|l\setminus F|-2<\theta$ is weak as well, because $\mu(\{l\})\le\sur_F(\{l\})$. For
  $\theta=1$ every remaining line is strong (\cref{prop:maxflow}\,(d)); in the default mode of the
  program (option \fn{--paranoid}, which the certifying run also gives explicitly) an assignment of excess $1$
  is constructed and validated for each of them nevertheless. For $\theta\ge 2$ a single-line flow with
  target $\theta$ is computed for each remaining line. Every line found strong is deleted at once.
\item \emph{Pairs.} If fewer than three lines are weak, the verdict is true. Otherwise a flow is computed
  for every $F$-adjacent pair $\{a,b\}$ of weak lines with $|(a\cup b)\setminus F|-4\ge\theta$; the other
  pairs are treated as not strong, which is always safe.
\item \emph{Triples.} A flow is computed for every $F$-connected triple of weak lines none of whose
  $F$-adjacent pairs was found strong; a pairwise $F$-adjacent triple is tested once. The verdict is
  false if one of these flows has gain less than $\theta$, and true otherwise.
\end{enumerate}
Every flow that is used for a true verdict is validated literally, in the form relative to the deleted
lines stated after \cref{lem:deleted}. A false verdict is reported together with the set of lines reached by the failed
search, whose surplus is recomputed from the definition. \fn{verify/spotcheck.py} implements
\cref{alg:leaf} with both pruning rules and without deletion; for pairs and triples it augments from one
line after the other, which gives the same gain by \cref{prop:maxflow}\,(e).

\subsubsection*{Cost}
Let $n=|L|$; at the leaves of the certificate $n$ is about $100$. One search for an augmenting path
inspects each live line at most once and costs $O(n)$ operations. The pairing needs $2n$ searches and the
single-line gains at most $n\theta$. A line is $F$-adjacent to at most $36$ others, so there are at most
$n\binom{36}{2}$ candidate triples, but in the certificate almost all of them contain a strong line. The
counters of the main verifier illustrate this. They are statistics of this one implementation; a verifier
written from \cref{alg:leaf} will not reproduce them unless it follows the steps (1)--(4) above.
Over the $9{,}179$ leaves with $\theta\ge2$ the main verifier computes $869{,}185$ single-line flows
(step~(2): one for every live line with at least $\theta+2$ cells outside $F$; from all these lines an
augmenting path exists), $77{,}268$ flows for pairs
(step~(3)) and only $7{,}391$ flows for triples (step~(4)). For the $32{,}703$ leaves with $\theta=1$ one
augmentation per line suffices, and these are not included in the counts. In all, the run makes
$8{,}549{,}521$ searches for augmenting paths (\fn{results/final/treecheck.json},
\fn{logs/repro/tc_all.log}). The three flow counts were recounted from the rules of steps (2)--(4) with
the matching code of the replication verifier of \cref{tab:verifiers}, without deletion of lines, and
came out the same (\fn{logs/paper/flow_counts/}); by \cref{lem:deleted}\,(b) the deletion does not change
which lines, pairs and triples are strong. For comparison, \cref{alg:leaf} exactly as printed, with pruning by strong
lines only, computes $4{,}063{,}331$ single-line gains (one for every live line at each of the $41{,}882$
leaves with $\theta\ge1$) and tests $692{,}694$ triples, each of them once (\fn{logs/paper/referee_replication.md}). In our runs a leaf test takes about $0.1$ milliseconds in the C++
verifier (with the literal validation of all flows) and about $2$ milliseconds in the Python verifier
($211$ CPU-seconds for all leaves). For $60.8\%$ of the leaves $\theta=0$ and only the pairing is
computed.

\subsection{Audits of the main verifier}\label{ssec:supp-audits}

This subsection gives the details of the three kinds of experiments of the audit of
\fn{verify/treecheck.cpp} (\fn{logs/review/treecheck_audit.md}) that \cref{ssec:audits} summarizes.

\subsubsection*{Mutations of the certificate}
In all, $66$ seeded mutations of the two files or of the options were submitted to the verifier. Of
these, $58$ were rejected: every deleted table entry (one for each round $1,\dots,10$), every replaced table move ($10$, half
of them busy cells), every move placed on a stone ($4$), every deleted leaf ($6$), every bogus leaf ($5$,
among them the root itself), Q-states listed as leaves ($2$), the bounds $\Rmax=9$ and $\Rmax=8$,
duplicate keys with a smaller move ($3$), an entry rewritten in another frame without mapping its move, all
truncated, empty, swapped or unreadable files ($9$), several damaged entries ($5$), and a run interrupted
by a time limit. Eight were accepted: the weaker bound $\Rmax=11$; two duplicate keys with a larger move, of which the verifier uses the smaller,
original one; entries rewritten in another frame \emph{with} the move mapped accordingly (three single
entries, and all entries at once), which the verifier re-canonicalizes; and additional well-formed entries for
unreachable states. None of the eight removes or changes an obligation: in each of them the visited
states and all counters are those of the unmodified run.

\subsubsection*{Mutations of the verifier}
Twelve variants of the source with one seeded defect each (danger children omitted; child (ii) omitted or
generated under a wrong condition; a listed leaf accepted without recomputation; wrong stone bounds;
threats on face diagonals or on space diagonals ignored; $\LOST$ always false; two defects in the
enumeration of triples in the leaf test; a too generous \VOID{} test) were compiled and run. The
self-test of the verifier detects all $12$. On the real certificate one mutant changes the verdict and seven change the
counters of \cref{tab:tree}; four change neither: the constant-false $\LOST$, which is harmless for
$\Rmax=10$, and the three defects that concern the leaf test (listed leaves not recomputed, and the two
defects in the enumeration of triples). The consequence of this last finding for the leaf test is drawn
in \cref{ssec:audits}.

\subsubsection*{Comparisons with brute force}
The leaf test of \fn{treecheck} was compared with a direct evaluation of the definition (Hall's condition
and the surplus condition over \emph{all} sets of live lines, for instances with at most $20$ live
lines): $53{,}755$ evaluations, including $5{,}970$ instances with $\theta$ exactly one above the true
minimum surplus, without a mismatch. The self-test of \fn{spotcheck.py} makes the same comparison on $6{,}000$
instances. The two leaf tests of the archive that are written in C++ (\fn{treecheck} and the library used
by the search) agree,
in verdict and in the reason for a failure, on all leaves and on $20{,}000$ perturbed instances derived
from the leaves, of which $13{,}688$ satisfy $\LEAF$ and $590$, $4{,}135$ and $1{,}587$ fail
conditions (L1), (L2) and (L3), respectively. $\LOST$ was compared with a literal transcription of the
patterns $\Etwo$ and $\Ethree$ on $800$ random admissible states with $t=9,\dots,13$, again without a mismatch.
During the internal refereeing, the audit referee's implementation of \cref{alg:leaf} was compared with
\fn{treecheck}, in verdict and in the failing condition, on the same $20{,}000$ perturbed instances, on all
leaves, on the leaf table and on the cross-check instances of \cref{ssec:supp-cross} with $k\ge|F|$, and the
replication verifier was compared with a direct evaluation of \cref{def:leaf} over all sets of live lines on
$19{,}995$ random instances with at most $16$ live lines; neither comparison found a mismatch
(\fn{logs/paper/referee_audit.md}, \fn{referee_replication.md}).

\subsection{Cross-checks that are not part of the proof}\label{ssec:supp-cross}

The following computations test the certificate and the programs from other directions. None of them is
needed for \cref{thm:main}; their coverage is stated exactly.

\subsubsection*{Exact enumeration of the pairing condition}
$\LEAF$ is only sufficient for the pairing condition $\CC(B,F,k)$ (\cref{thm:leaf-implies-C}). For small
$\theta$ the statement $\CC(B,F,t)$ itself can be decided by brute force: the program \fn{bin/enumcheck}
enumerates, depth first and in ascending order of the cells, all $B$-admissible sets $X\supseteq F$ with
$|X|\le t$, up to two reductions that do not change the answer (cells on no live line are skipped, and
a set whose last stone lies on a cell that is unassigned in the current pairing is not examined); it
maintains a pairing by augmenting paths and stops when some $X$ has none. It was run on all
leaves with $\theta\le3$, that is on $64{,}841+32{,}703+7{,}267+1{,}434=106{,}245$ leaves ($99.55\%$): the
pairing condition holds for every one of them (\fn{results/final/checks/enum_theta*.jsonl}). The
largest single enumeration examines $1.39$ million sets. The programs \fn{bin/enumcheck} and
\fn{bin/certify} (below) are built on the library of the search (\fn{src/match.hpp}, \fn{src/phi.hpp});
they are cross-checks of the certificate and are not independent of the search program.

\subsubsection*{The integer program}
\Cref{sec:ip} describes a $0$--$1$ program whose optimal value is positive if and only if $\CC(B,F,k)$
fails. It was solved with the floating-point solver HiGHS~\cite{HuangfuHall2018}
\begin{itemize}[leftmargin=1.6em]
\item on all $478$ leaves with $\theta\ge4$ and on a seeded random sample of $830$ leaves with
  $\theta\le 3$: all $1{,}308$ answers are ``holds'';
\item on the leaf table of \cref{ssec:leaftable}: for $k=8$ on all $6$ orbits with slack $8$ (all hold); for
  $k=9$ on $262$ of the $947$ orbits with slack at least $9$ ($254$ fail, $8$ hold), in agreement with the
  max-flow table in all $268$ cases; the remaining orbits for $k=9$ and all orbits for $k=10$ were not run;
\item on $1{,}291$ random and structured instances, against an exact integer decision procedure for $\CC$
  (\fn{bin/certify}: the leaf test, an explicit admissible violator, exact enumeration, or a complete branch
  and bound). The exact procedure decides all of them: $656$ hold, $623$ fail, and $12$ are degenerate
  (in $9$ of them $F$ is not $B$-admissible, in $3$ of them $|F|>k$; the integer program is infeasible in
  both cases, \cref{prop:ip-exact}). The integer program decides $1{,}275$ of the $1{,}279$ non-degenerate
  instances ($652$ hold,
  $623$ fail) with no disagreement, reports the $12$ degenerate instances as infeasible, and runs into the time
  limit of $300$ seconds on $4$ instances, all with $F=\emptyset$ and all holding by $\LEAF$. Each of the
  $623$ violating positions found by one program was re-checked with the predicates of the other.
\end{itemize}
Among the $656$ instances that hold there are $160$ for which $\LEAF$ is false: $\LEAF$ is a strictly
stronger condition than $\CC$ (it fails on instances on which $\CC$ holds), not only in the situation
of \cref{ex:C-not-leaf}. No such state is needed by the certificate.

\subsubsection*{Simulated games}
The program \fn{sim/sim.cpp} ($4{,}383$ lines, written from the specification, sharing no code with the
other programs) plays the real strategy $\sigma$ of \cref{def:sigma} on the real board: it looks up the two
files, answers threats and danger pairs, reveals occupied cells, computes an actual pairing for the real
position when a leaf is reached, and plays the pairing strategy to the full board. It reports a completed
line of Maker, a missing table entry, a leaf whose real position has no pairing, a position that is not
admissible after Breaker's move, and an illegal table move. In $10^6$ games against a mixture of nine
adversaries (random; greedy line building; stealing the cells Breaker wants; making threats; a two-ply
search against $\sigma$; building danger pairs; two successive threats; occupying the corners of a face;
and an adversary that uses the table itself to steer into rarely visited states) none of these events
occurred: $0$ lines of Maker, $0$ missing entries, $0$ missing pairings, $0$ inadmissible positions and $0$
illegal moves in $10^6$ completed games with $1{,}121{,}484$ threat rounds, $85{,}806$ danger rounds and $2{,}254{,}365$ reveals
(\fn{results/final/sim/sim_mix_1M.json}). The leaf was reached after round $8$, $9$, $10$ in
$117{,}835$, $832{,}519$ and $49{,}646$ games.

This is a sampling test and nothing more. The games met $78{,}957$ of the $93{,}164$ table entries
($84.8\%$) and $81{,}907$ of the $106{,}723$ leaves ($76.7\%$); $14{,}207$ table entries and $24{,}816$
leaves, more than nine tenths of them in rounds $9$ and $10$, were never played. Those states are
covered by the verifiers of \cref{ssec:impl} only.

\subsubsection*{Summary for the leaves}
\Cref{tab:leaves} in \cref{ssec:cross} shows how each leaf is covered by the verifiers and by the cross-checks of
this subsection.

\subsection{Known minor issues of the tools}\label{ssec:supp-issues}

The audits found no defect that affects a result, but the following properties of the tools should be
known to anyone who reuses them. \fn{treecheck} accepts a table that contains two different moves for the
same key (it uses the smaller move and counts the conflict), accepts entries written in a non-canonical
frame (it re-canonicalizes them and counts them), leaves an older output file in place if it aborts,
parses numeric options leniently, and reports success for a \VOID{} root given on the command line; hence
the list of output fields to be checked in \cref{ssec:impl}. All these counters are zero in the certified
run. The statistics file of the search counts only \VOID{} children of types (c) and (d) (there are none),
not the $151$ \VOID{} children of type (ii). The $\LOST$ code of all programs is not exercised by a
certificate with $\Rmax=10$; it would become load-bearing for $\Rmax\ge11$. The simulator does not count a
missing table entry that is met only while an adversary predicts Breaker's next move.
